% Made up for the tikra gallery: an inscribed angle is half the central angle on the same arc.
% In the playground, drag A, B or C around the circle: the two angles keep their ratio.
\begin{tikzpicture}[scale=2]
  \def\a{200} % where A is on the circle, in degrees, 130..350
  \def\b{-30} % where B is, -60..20
  \def\c{60}  % where C is, 30..110
  \useasboundingbox (-2.3,-1.3) rectangle (1.45,1.3);
  \draw[thick] (0,0) circle (1);
  \coordinate (O) at (0,0);
  \coordinate (A) at (\a:1);
  \coordinate (B) at (\b:1);
  \coordinate (C) at (\c:1);
  \draw[thick, blue!75!black, fill=blue!8] (B) -- (A) -- (C);
  \draw[thick, red!75!black, fill=red!8] (B) -- (O) -- (C);
  \draw[ultra thick, violet] (\b:1) arc (\b:\c:1);
  \foreach \p/\w in {A/left, B/below right, C/above right} \fill (\p) circle (0.9pt) node[\w] {$\p$};
  \fill (O) circle (0.8pt) node[left] {$O$};
  \pgfmathsetmacro\central{\c-\b}
  \node[font=\small, align=left, anchor=north west] at (-2.28,1.28) {$\angle BOC = \pgfmathprintnumber[fixed, precision=0]{\central}^\circ$\\
    $\angle BAC = \pgfmathprintnumber[fixed, precision=1]{\central/2}^\circ$};
\end{tikzpicture}
